\chapter{Dynamic Reconfiguration and Intelligent Offloading}
\label{chap:dynamic-reconfiguration}

The macroprogramming models of \cref{chap:programming-models} abstract from physical placement by design:
they fix what each logical participant runs,
not which physical host runs it.
%
Some mechanism must therefore supply the missing choice of host at run time,
and keep revising it as the infrastructure described in \cref{chap:edge-cloud-continuum} changes underneath it.
%
This chapter reviews four answers to that need:
constraint- and optimisation-based placement,
energy- and carbon-aware objectives,
declarative statements of the deployment problem,
and policies learned with \ac{RL} over graph representations of the topology.
%
They are the background for the thesis' own answers in Part III:
the deployment and reconfiguration strategies of \cref{chap:deployment-and-reconfiguration} and the learned offloading policy of \cref{chap:heterogeneous-gnn}.

\section{Deployment Problems in Pervasive Systems}
\label{sec:the-deployment-problem}

Deployment is a challenge shared by the target environments of this thesis (\cref{subsec:environment-challenges}):
it means placing computation and coordination logic onto nodes that are heterogeneous (\cref{subsec:continuum-characteristics}) and mobile or intermittently connected (\cref{subsec:network-volatility}),
so that no placement decision can be assumed to remain valid for long.
%
Once the assumption of a fixed topology is dropped,
the mapping from logical computational units to physical hosts has to be decided,
revised,
and enacted at run time.

Deployment is an assignment problem:
given a set of computational units -- services, functions, actors -- and a set of candidate hosts connected by links of varying capacity,
find a mapping from the former to the latter that respects the units' resource and connectivity requirements (\cref{fig:placement-mapping}).
%
\emph{Static deployment} answers this question once, at design or release time, and keeps the resulting mapping fixed for the lifetime of the application;
this is still the default assumption behind most automatic deployment tooling, which targets getting a system correctly installed and running rather than keeping it correctly placed as conditions change~\cite{DBLP:journals/jss/ArcangeliBL15}.
%
\emph{Dynamic deployment} instead treats the mapping as a variable that must be revisited continuously,
recomputing or adjusting it as devices join,
leave,
move,
or change state -- the behaviour argued necessary for self-managed systems in general in \cref{subsec:continuum-evolution}~\cite{DBLP:conf/icse/KramerM07}.

\begin{figure}[t]
    \centering
    % Placement as a mapping between two graphs (application -> infrastructure).
% Requires in the preamble: \usepackage{tikz}
%   \usetikzlibrary{arrows.meta,positioning,fit,backgrounds,shapes.geometric}
% Include with a normal figure environment, e.g. % Placement as a mapping between two graphs (application -> infrastructure).
% Requires in the preamble: \usepackage{tikz}
%   \usetikzlibrary{arrows.meta,positioning,fit,backgrounds,shapes.geometric}
% Include with a normal figure environment, e.g. % Placement as a mapping between two graphs (application -> infrastructure).
% Requires in the preamble: \usepackage{tikz}
%   \usetikzlibrary{arrows.meta,positioning,fit,backgrounds,shapes.geometric}
% Include with a normal figure environment, e.g. \input{figures/placement-mapping}.
\begin{tikzpicture}[
  >=Latex,
  font=\small,
  comp/.style={circle,draw,minimum size=6mm,inner sep=0pt,fill=gray!12},
  host/.style={rounded corners,draw,align=center,inner sep=2pt},
  cloud/.style={host,fill=gray!8,minimum width=18mm,minimum height=9mm},
  edge/.style={host,fill=gray!16,minimum width=15mm,minimum height=8mm},
  gw/.style={host,fill=gray!24,minimum width=13mm,minimum height=7mm},
  dev/.style={host,fill=gray!34,minimum width=11mm,minimum height=6mm},
  link/.style={draw,thick},
  clink/.style={draw,gray!60},
  map/.style={->,dashed,gray!70,shorten >=1pt,shorten <=1pt},
]

% ---------- application graph (left) ----------
\node[comp] (c1) at (0,2.2) {$c_1$};
\node[comp] (c2) at (1.4,2.9) {$c_2$};
\node[comp] (c3) at (1.5,1.3) {$c_3$};
\node[comp] (c4) at (0.2,0.5) {$c_4$};

\draw[clink] (c1) -- (c2);
\draw[clink] (c1) -- (c3);
\draw[clink] (c3) -- (c4);
\draw[clink] (c2) -- (c3);

\begin{scope}[on background layer]
  \node[draw,dotted,rounded corners,fit=(c1)(c2)(c3)(c4),inner sep=6pt] (app) {};
\end{scope}
\node[anchor=south] at (app.north) {\textbf{Application graph}};
\node[anchor=north,font=\scriptsize,text width=3.2cm,align=center]
  at (app.south) {components + communication demands};

% ---------- infrastructure graph (right) ----------
\node[cloud] (cl) at (7.2,3.0) {Cloud};
\node[edge]  (e1) at (6.3,1.3) {Edge};
\node[gw]    (g1) at (8.3,1.4) {Gateway};
\node[dev]   (d1) at (5.9,-0.3) {Device};
\node[dev]   (d2) at (8.6,-0.2) {Robot};

\draw[link] (cl) -- (e1);
\draw[link] (cl) -- (g1);
\draw[link] (e1) -- (d1);
\draw[link] (e1) -- (g1);
\draw[link] (g1) -- (d2);

\begin{scope}[on background layer]
  \node[draw,dotted,rounded corners,fit=(cl)(e1)(g1)(d1)(d2),inner sep=6pt] (infra) {};
\end{scope}
\node[anchor=south] at (infra.north) {\textbf{Infrastructure graph}};
\node[anchor=north,font=\scriptsize,text width=3.6cm,align=center]
  at (infra.south) {heterogeneous hosts + links (latency, bandwidth)};

% ---------- mapping ----------
\draw[map] (c2) to[out=20,in=160] (cl);
\draw[map] (c1) to[out=-10,in=170] (e1);
\draw[map] (c3) to[out=0,in=220] (g1);
\draw[map] (c4) to[out=-20,in=190] (d1);

\node[font=\scriptsize\itshape,gray!70,fill=white,inner sep=1.5pt] at (3.9,3.2) {placement (mapping)};

\end{tikzpicture}
.
\begin{tikzpicture}[
  >=Latex,
  font=\small,
  comp/.style={circle,draw,minimum size=6mm,inner sep=0pt,fill=gray!12},
  host/.style={rounded corners,draw,align=center,inner sep=2pt},
  cloud/.style={host,fill=gray!8,minimum width=18mm,minimum height=9mm},
  edge/.style={host,fill=gray!16,minimum width=15mm,minimum height=8mm},
  gw/.style={host,fill=gray!24,minimum width=13mm,minimum height=7mm},
  dev/.style={host,fill=gray!34,minimum width=11mm,minimum height=6mm},
  link/.style={draw,thick},
  clink/.style={draw,gray!60},
  map/.style={->,dashed,gray!70,shorten >=1pt,shorten <=1pt},
]

% ---------- application graph (left) ----------
\node[comp] (c1) at (0,2.2) {$c_1$};
\node[comp] (c2) at (1.4,2.9) {$c_2$};
\node[comp] (c3) at (1.5,1.3) {$c_3$};
\node[comp] (c4) at (0.2,0.5) {$c_4$};

\draw[clink] (c1) -- (c2);
\draw[clink] (c1) -- (c3);
\draw[clink] (c3) -- (c4);
\draw[clink] (c2) -- (c3);

\begin{scope}[on background layer]
  \node[draw,dotted,rounded corners,fit=(c1)(c2)(c3)(c4),inner sep=6pt] (app) {};
\end{scope}
\node[anchor=south] at (app.north) {\textbf{Application graph}};
\node[anchor=north,font=\scriptsize,text width=3.2cm,align=center]
  at (app.south) {components + communication demands};

% ---------- infrastructure graph (right) ----------
\node[cloud] (cl) at (7.2,3.0) {Cloud};
\node[edge]  (e1) at (6.3,1.3) {Edge};
\node[gw]    (g1) at (8.3,1.4) {Gateway};
\node[dev]   (d1) at (5.9,-0.3) {Device};
\node[dev]   (d2) at (8.6,-0.2) {Robot};

\draw[link] (cl) -- (e1);
\draw[link] (cl) -- (g1);
\draw[link] (e1) -- (d1);
\draw[link] (e1) -- (g1);
\draw[link] (g1) -- (d2);

\begin{scope}[on background layer]
  \node[draw,dotted,rounded corners,fit=(cl)(e1)(g1)(d1)(d2),inner sep=6pt] (infra) {};
\end{scope}
\node[anchor=south] at (infra.north) {\textbf{Infrastructure graph}};
\node[anchor=north,font=\scriptsize,text width=3.6cm,align=center]
  at (infra.south) {heterogeneous hosts + links (latency, bandwidth)};

% ---------- mapping ----------
\draw[map] (c2) to[out=20,in=160] (cl);
\draw[map] (c1) to[out=-10,in=170] (e1);
\draw[map] (c3) to[out=0,in=220] (g1);
\draw[map] (c4) to[out=-20,in=190] (d1);

\node[font=\scriptsize\itshape,gray!70,fill=white,inner sep=1.5pt] at (3.9,3.2) {placement (mapping)};

\end{tikzpicture}
.
\begin{tikzpicture}[
  >=Latex,
  font=\small,
  comp/.style={circle,draw,minimum size=6mm,inner sep=0pt,fill=gray!12},
  host/.style={rounded corners,draw,align=center,inner sep=2pt},
  cloud/.style={host,fill=gray!8,minimum width=18mm,minimum height=9mm},
  edge/.style={host,fill=gray!16,minimum width=15mm,minimum height=8mm},
  gw/.style={host,fill=gray!24,minimum width=13mm,minimum height=7mm},
  dev/.style={host,fill=gray!34,minimum width=11mm,minimum height=6mm},
  link/.style={draw,thick},
  clink/.style={draw,gray!60},
  map/.style={->,dashed,gray!70,shorten >=1pt,shorten <=1pt},
]

% ---------- application graph (left) ----------
\node[comp] (c1) at (0,2.2) {$c_1$};
\node[comp] (c2) at (1.4,2.9) {$c_2$};
\node[comp] (c3) at (1.5,1.3) {$c_3$};
\node[comp] (c4) at (0.2,0.5) {$c_4$};

\draw[clink] (c1) -- (c2);
\draw[clink] (c1) -- (c3);
\draw[clink] (c3) -- (c4);
\draw[clink] (c2) -- (c3);

\begin{scope}[on background layer]
  \node[draw,dotted,rounded corners,fit=(c1)(c2)(c3)(c4),inner sep=6pt] (app) {};
\end{scope}
\node[anchor=south] at (app.north) {\textbf{Application graph}};
\node[anchor=north,font=\scriptsize,text width=3.2cm,align=center]
  at (app.south) {components + communication demands};

% ---------- infrastructure graph (right) ----------
\node[cloud] (cl) at (7.2,3.0) {Cloud};
\node[edge]  (e1) at (6.3,1.3) {Edge};
\node[gw]    (g1) at (8.3,1.4) {Gateway};
\node[dev]   (d1) at (5.9,-0.3) {Device};
\node[dev]   (d2) at (8.6,-0.2) {Robot};

\draw[link] (cl) -- (e1);
\draw[link] (cl) -- (g1);
\draw[link] (e1) -- (d1);
\draw[link] (e1) -- (g1);
\draw[link] (g1) -- (d2);

\begin{scope}[on background layer]
  \node[draw,dotted,rounded corners,fit=(cl)(e1)(g1)(d1)(d2),inner sep=6pt] (infra) {};
\end{scope}
\node[anchor=south] at (infra.north) {\textbf{Infrastructure graph}};
\node[anchor=north,font=\scriptsize,text width=3.6cm,align=center]
  at (infra.south) {heterogeneous hosts + links (latency, bandwidth)};

% ---------- mapping ----------
\draw[map] (c2) to[out=20,in=160] (cl);
\draw[map] (c1) to[out=-10,in=170] (e1);
\draw[map] (c3) to[out=0,in=220] (g1);
\draw[map] (c4) to[out=-20,in=190] (d1);

\node[font=\scriptsize\itshape,gray!70,fill=white,inner sep=1.5pt] at (3.9,3.2) {placement (mapping)};

\end{tikzpicture}

    \caption[Deployment as an assignment problem between two graphs]{Deployment as an assignment problem: an application graph of components annotated with their communication demands (left) is mapped onto an infrastructure graph of heterogeneous hosts and capacity-, latency-, and bandwidth-annotated links (right). Static deployment fixes this mapping once; dynamic deployment revisits it as hosts join, leave, move, or change state.}
    \label{fig:placement-mapping}
\end{figure}

Static deployment is the wrong model once the assumptions that justify it stop holding.
%
A mapping computed once assumes that the set of candidate hosts,
their capacity,
and the connectivity between them are known and stable enough for a single decision to remain valid.
%
None of the three holds towards the edge of the continuum (\cref{subsec:continuum-characteristics,subsec:network-volatility}),
where nodes differ widely in capacity and energy budget and links form and break as devices move.
%
Under these conditions a static mapping does not merely become suboptimal over time.
%
It can stop being satisfiable if the host it names disappears or the link to it is no longer available.

Reopening the mapping at run time raises four questions, and they structure the rest of this chapter:
which objective a placement should optimise, and under which constraints (\cref{subsec:constraint-based-placement});
which objective to privilege when several compete, energy being one of them (\cref{subsec:energy-aware-orchestration});
in what form the components, their requirements, and the objective are stated to whatever computes the mapping (\cref{subsec:declarative-deployment});
and how a new mapping can be computed fast enough to matter on a topology that keeps changing while the computation runs (\cref{sec:ml-for-topologies}).
%
The last question drives the shift from the classical heuristics surveyed in \cref{sec:edge-orchestration-sota} to learning-based methods.

\section{State-of-the-Art in Edge Orchestration and Placement}
\label{sec:edge-orchestration-sota}

Automating the mapping problem posed in \cref{sec:the-deployment-problem} is the subject of a substantial body of work on \emph{edge orchestration and service placement}.
%
Aït-Salaht et al. survey it, proposing a classification scheme for the field and cataloguing its open issues~\cite{DBLP:journals/csur/Ait-SalahtDL20}.
%
Taleb et al. survey it again four years later, which is itself a sign that the first wave of proposals left placement in integrated cloud-fog-edge settings unsettled~\cite{DBLP:journals/csur/TalebGD25}.
%
The later survey groups existing proposals into three categories:
graph-based (exact) methods,
heuristics,
and machine-learning-based methods~\cite{DBLP:journals/csur/TalebGD25}.
%
This chapter follows that classification loosely, with one addition.
%
The first two categories are treated together under constrained optimisation, with energy as an objective of its own (\cref{subsec:constraint-based-placement,subsec:energy-aware-orchestration}).
%
The addition is how the problem itself is stated, rather than how it is solved, a concern the classification leaves implicit (\cref{subsec:declarative-deployment}).
%
The third category, learning-based methods, is the most active recent trend in both surveys (\cref{sec:ml-for-topologies,fig:orchestration-taxonomy}).

\begin{figure}[t]
    \centering
    \resizebox{\textwidth}{!}{% Taxonomy of edge-orchestration and placement methods (after Taleb et al.).
% Requires: \usepackage[edges]{forest}
% Wrap in \resizebox{\textwidth}{!}{...} in the figure environment: this tree is
% naturally wider than \textwidth.
\begin{forest}
  for tree={
    draw, rounded corners, align=center, font=\small,
    parent anchor=south, child anchor=north,
    edge={-,gray}, l sep=8mm, s sep=6mm, inner sep=3pt,
  }
  [Edge orchestration \& placement, fill=gray!8
    [Graph-based\\(exact), fill=gray!16
      [{ILP / CP over two graphs\\\scriptsize Skarlat et al.; Brogi \& Forti;\\\scriptsize Miloradovi\'{c} \& Papadopoulos}, fill=white]
    ]
    [Heuristic /\\metaheuristic, fill=gray!16
      [{consolidation, semi-greedy, GA\\\scriptsize Beloglazov et al.; Azizi et al.}, fill=white]
    ]
    [Machine\\learning, fill=gray!16
      [{online bandit; deep RL + GNN\\\scriptsize Struh\'{a}r et al.; Huang et al.}, fill=white]
    ]
  ]
\end{forest}
}
    \caption[Taxonomy of edge-orchestration and placement methods]{The three categories into which Taleb et al.~\cite{DBLP:journals/csur/TalebGD25} group edge-orchestration and placement proposals, with representative works from this chapter as leaves. This chapter reviews the first two under constrained optimisation (\cref{subsec:constraint-based-placement,subsec:energy-aware-orchestration}) and turns to the third, learning-based, category in \cref{sec:ml-for-topologies}.}
    \label{fig:orchestration-taxonomy}
\end{figure}

\subsection{Optimisation and Constraint-based Placement}
\label{subsec:constraint-based-placement}

The dominant formulation treats placement as a constrained combinatorial optimisation problem over two graphs.
%
An application is described as a graph of components annotated with their resource demands and mutual communication requirements,
and the infrastructure as a graph of nodes annotated with capacity and links annotated with latency and bandwidth.
%
A placement is then an assignment of the first graph onto the second that respects every node's and link's capacity while optimising an objective function~\cite{DBLP:journals/csur/Ait-SalahtDL20}.
%
Proposals in what Taleb et al. call the graph-based category typically pose this assignment as an \ac{ILP} over the two graphs, solvable exactly for small instances~\cite{DBLP:journals/csur/TalebGD25}.
%
Skarlat et al. instantiate the general formulation for \ac{IoT} services over fog resources, explicitly weighing the \ac{QoS} attributes of both applications and resources in the objective, and pair an exact solver with a genetic-algorithm heuristic once instance sizes make the exact method impractical~\cite{DBLP:journals/soca/SkarlatNSBL17}.
%
Brogi and Forti tackle the same assignment from the eligibility side: their model for multi-component \ac{IoT} applications over cloud-fog infrastructure, prototyped in the FogTorch tool, filters candidate placements by the hardware, software, and \ac{QoS} constraints each component and host must jointly satisfy, rather than searching for a single numerically optimal one~\cite{DBLP:journals/iotj/BrogiF17}.
%
Miloradovic and Papadopoulos solve the same assignment for application offloading over what they term the Edge-to-Cloud Continuum in two ways, formulating it once as an \ac{ILP} and once as a constraint-programming model~\cite{DBLP:conf/cdc/MiloradovicP23}.
%
Their objective folds the carbon footprint of the offloaded computation in alongside resource utilisation.
%
The same pattern recurs across these proposals, and across the graph-based category more broadly.
%
Exact methods give an optimal answer for a snapshot of the infrastructure, but their cost grows too fast with the number of nodes and services to be recomputed as often as a volatile continuum demands (\cref{subsec:network-volatility}).
%
Heuristics -- and, later, the learning-based methods of \cref{sec:ml-for-topologies} -- therefore take over as problem size and reconfiguration frequency both grow.

A distinct line of work brings the reconfiguration-frequency concern into the container runtime itself, for applications with real-time constraints that ordinary container orchestration does not preserve~\cite{DBLP:conf/fogiot/StruharBAP20}.
%
Struhár et al. decompose the problem into two phases.
%
An offline phase computes an optimal initial placement and dimensioning of containers from a system model.
%
An online phase then adapts each container's resource allocation at run time in response to workload variation and interference,
without which a real-time containerised application would drift out of its timing requirements~\cite{DBLP:journals/tecs/StruharCABP24}.
%
A companion proposal narrows the online objective, replacing the generic timing bound with a specific one: it adapts a containerised control application's resource allocation to bound the deviation of its response from the ideal~\cite{DBLP:conf/etfa/StruharABPC23}.
%
Both examples make that split explicit: reconfiguration is part of the deployment model, a second phase running continuously alongside the initial placement.

Operational practice lags behind this literature.
%
Cilic et al. benchmark existing edge orchestration platforms against the \ac{QoS} attributes they target and their ability to schedule work onto remote, edge-located resources~\cite{DBLP:journals/sensors/CilicKZK23}.
%
They find Kubernetes and its distributions the most ready deployment vehicle available today, but still short of the dynamic, distributed character of edge environments.
%
This gap matches the limitation identified in the underlying architectural paradigms (\cref{subsec:paradigm-limitations}):
the placement policy belongs to the orchestrator,
stays comparatively coarse,
and is never exposed to the application as a first-class, programmable artefact able to host the richer objective functions surveyed above.

\subsection{Energy-aware Orchestration}
\label{subsec:energy-aware-orchestration}

Energy is an axis of heterogeneity that does not track computational capacity (\cref{subsec:continuum-characteristics}), and it is a deployment concern in its own right.
%
Energy-aware placement predates the edge.
%
In the data centre, Beloglazov et al. cast it as a consolidation problem: pack virtual machines onto the fewest active hosts that still meet each workload's negotiated \ac{QoS}, and power down or idle the rest~\cite{DBLP:journals/fgcs/BeloglazovAB12}.
%
Because no exact solver handles that scale, they rely on heuristics aimed at the energy wasted by idle servers.
%
That formulation assumes hosts of roughly uniform energy cost, with consolidation as the main lever.
%
Neither assumption survives the move to the edge, where energy budgets vary by orders of magnitude between nodes (\cref{subsec:continuum-characteristics}), from mains-powered gateways to battery-limited robots to intermittent harvested supplies.
%
At that end the supply is also unreliable:
an energy-harvesting device can lose power in the middle of a computation and resume only once it has stored enough charge to continue,
so that its computation advances in bursts separated by power failures~\cite{DBLP:conf/snapl/LuciaBCMR17}.
%
Orchestration cannot then treat such a node as continuously available for the work assigned to it, which rules out a placement decided once and left to stand.

Fog- and edge-specific proposals adjust the objective accordingly.
%
Azizi et al., for instance, schedule \ac{IoT} tasks across fog nodes with semi-greedy heuristics that jointly minimise total energy consumption and deadline-violation time, in place of data-centre-style consolidation~\cite{DBLP:journals/jnca/AziziSAB22}.
%
Their scheduler is one of many heuristic and metaheuristic fog schedulers that pursue the same energy/deadline trade-off under different search strategies.
%
These schedulers still treat energy as a quantity to be minimised wherever and whenever the work runs.
%
Wiesner et al. change what is counted, shifting the objective from energy to carbon.
%
Emissions per unit of electricity vary through the day as the grid's generation mix changes, so they defer flexible workloads to the intervals when the supply is least carbon-intensive, guided by forecasts of that intensity~\cite{DBLP:conf/middleware/WiesnerBSGT21}.
%
A placement should then track when it draws energy, not only how much.
%
At the edge, a node's energy budget is not fixed either: when it is drawn from renewable or harvested sources rather than a stable grid, the available supply fluctuates over time.
%
Struhár et al. address this by framing container placement at the edge as a multi-armed bandit problem~\cite{DBLP:conf/ucami/StruharPAAF25}.
%
The orchestrator then learns online which placement to prefer as each node's green-energy supply fluctuates, instead of optimising a fixed budget assumed at deployment time.
%
How such an objective is stated to a planner, together with the components and the infrastructure, is a separate question, taken up next.

\subsection{Declarative Deployment Specification}
\label{subsec:declarative-deployment}

The formulations reviewed in \cref{subsec:constraint-based-placement,subsec:energy-aware-orchestration} presuppose that the problem is already stated: which components exist, what they require, what the infrastructure offers, and which objective the mapping should optimise.
%
Industrial practice states this declaratively.
%
TOSCA, the \emph{Topology and Orchestration Specification for Cloud Applications}, standardises a portable model of an application's components, their relationships, and the operations that manage them, so that a deployment can be described once and enacted by different orchestrators~\cite{DBLP:books/sp/aws14/BinzBKL14}.
%
\emph{Infrastructure as code} generalises the practice: provisioning and configuration are written as versioned, executable specifications rather than performed by hand~\cite{DBLP:journals/infsof/RahmanMW19}.
%
Kubernetes' operating model is itself declarative: the operator declares a desired state, and controllers continuously reconcile the running system towards it.
%
The declared state, however, remains an operations artefact: it belongs to whoever runs the platform, not to the application -- the same boundary found around the placement policy itself in \cref{subsec:constraint-based-placement}.

The declarative statement can also drive the placement search.
%
The constraint-filtering model of Brogi and Forti, reviewed in \cref{subsec:constraint-based-placement}, develops into \emph{declarative continuous reasoning}: the deployment goal is stated once, and an incremental reasoner repairs the placement as the infrastructure drifts, keeping the goal satisfied instead of re-solving from scratch at every change~\cite{DBLP:journals/logcom/FortiBB22}.
%
For collective systems, the declarative statement need not stop at operations: it can become the artefact the application itself manipulates at run time, the direction taken for pulverised deployments~\cite{DBLP:journals/fgcs/FarabegoliPCV24}.

Two of these strands are combined into a declarative planner for green deployments on the continuum in \cref{chap:deployment-and-reconfiguration}:
energy and carbon objectives tracked against a time-varying supply (\cref{subsec:energy-aware-orchestration}),
and a deployment goal stated once and kept satisfied.
%
Placement can also be learned online instead of declared and repaired, which is the subject of \cref{sec:ml-for-topologies}.

\section{Machine Learning for Distributed Network Topologies}
\label{sec:ml-for-topologies}

The methods reviewed in \cref{sec:edge-orchestration-sota} share a weakness under the conditions established in \cref{subsec:network-volatility}:
each recomputes a placement or a schedule from a model of the infrastructure assumed accurate at the moment the computation runs.
%
A mobile, intermittently connected continuum keeps violating that assumption.
%
Struhár et al.'s bandit formulation (\cref{subsec:energy-aware-orchestration}) is the first proposal reviewed here to loosen it, learning online which placement to prefer instead of re-solving against a budget assumed fixed at deployment.
%
Machine learning generalises that move, and both surveys cited in \cref{sec:edge-orchestration-sota} report it as the most active recent trend in edge orchestration~\cite{DBLP:journals/csur/Ait-SalahtDL20,DBLP:journals/csur/TalebGD25}.

A learned policy takes a different route to the same decision.
%
Instead of solving an optimisation problem anew against the latest known state, it is trained to map states of the environment directly to placement or offloading decisions, whatever form those states take.
%
This moves the cost of search into a training phase that runs once, offline; at run time the decision reduces to a single forward pass through the learned function.
%
The methods examined below separate training from operation more sharply than the online bandit, which learns as it runs: they train against a model of the environment ahead of time, before any decision is taken.
%
This separation pays off only if the trained policy generalises to states it did not meet during training, a requirement returned to in \cref{subsec:gnns-for-edge}.
%
Not every kind of learning fits this setting.
%
Supervised learning would need a supply of ground-truth optimal placements to imitate, the expensive optimisation outputs the learned approach exists to avoid computing.
%
\ac{RL}, by contrast, asks only for a reward signal that scores a placement after the fact.
%
A survey of deep \ac{RL} in communications and networking catalogues its application across this class of resource-allocation and offloading problems~\cite{DBLP:journals/comsur/LuongHGNWLK19}, and it is the route the rest of this section, and \cref{chap:heterogeneous-gnn}, take.

The three subsections that follow take that route in order:
task offloading, where the decision narrows to a single unit and a binary choice of host, so the contrast with a re-solved optimisation is sharpest (\cref{subsec:learning-based-offloading});
the \ac{RL} formulation and the deep Q-learning algorithm that realises it at scale (\cref{subsec:reinforcement-learning});
and the \acp{GNN} needed to feed such a policy a representation of a topology that is itself heterogeneous and changing (\cref{subsec:gnns-for-edge}).
%
The combination is later developed into a full offloading policy (\cref{chap:heterogeneous-gnn}).

\subsection{Traditional vs. Learning-based Offloading}
\label{subsec:learning-based-offloading}

Task offloading is the decision, taken by or on behalf of a resource-constrained device, of whether to execute a computation locally or to send it to a more capable node reachable over the network~\cite{DBLP:journals/csur/Ait-SalahtDL20}.
%
Cast this way, offloading is a special case of the placement problem of \cref{sec:the-deployment-problem}, restricted to a single computational unit and a binary choice of host.
%
The exact and heuristic methods of \cref{subsec:constraint-based-placement} therefore apply to it unchanged: an offloading policy can be obtained by solving an \ac{ILP} or a constraint-satisfaction problem over the device's and the remote host's resources, recomputed whenever the channel condition or the workload changes.
%
Huang et al. take the same problem to a learned policy instead, training an agent to decide, for each arriving task, whether to compute locally or offload over a wireless-powered link, and to do so directly from the time-varying channel state rather than from a re-solved optimisation problem~\cite{DBLP:journals/tmc/HuangBZ20}.
%
Their network learns these binary decisions from experience, dispensing with the combinatorial search the constraint-based formulation would rerun for every new channel state.
%
The learned policy does not reach a different answer from the heuristic on a fixed instance:
on a static snapshot of the channel, solving the problem exactly still works.
%
The two differ in where the cost of adapting to change is paid.
%
An optimisation-based policy pays that cost at every decision, re-solving the problem instance whenever the channel or workload changes.
%
A learned policy pays it once, during training, and reduces every later decision to a forward pass through the trained model (\cref{fig:cost-shift}).
%
That matters under the conditions already established for the continuum (\cref{subsec:network-volatility}), where the channel and the topology change faster than a fresh optimisation could be solved and applied.
%
It does not remove the need for a model of the problem: the agent still needs a representation of states, actions, and rewards, developed next in \cref{subsec:reinforcement-learning}.
%
The burden moves from finding a solution online to finding one during a training phase that need not run in real time.

\begin{figure}[t]
    \centering
    % Where the adaptation cost is paid: re-solved optimisation vs. learned policy.
% Requires: \usepackage{tikz}
%   \usetikzlibrary{arrows.meta,positioning,fit,backgrounds,decorations.pathreplacing}
\begin{tikzpicture}[
  >=Latex, font=\small,
  solve/.style={draw,fill=gray!32,minimum height=9mm,minimum width=13mm,align=center,inner sep=1pt},
  fpass/.style={draw,fill=gray!12,minimum height=6mm,minimum width=8mm,align=center,inner sep=1pt},
  train/.style={draw,fill=gray!22,minimum height=15mm,minimum width=20mm,align=center},
  dec/.style={circle,fill=black,inner sep=1.4pt},
  ax/.style={->,gray!55,thick},
]
% offline | online separator
\draw[dotted,gray] (2.55,-1.2) -- (2.55,3.1);
\node[gray,font=\scriptsize,anchor=south] at (1.35,3.55) {offline (once)};
\node[gray,font=\scriptsize,anchor=south] at (7.0,3.55) {online (per decision)};

% lane labels
\node[anchor=east,align=right] at (-0.15,2.1) {\textbf{Re-solved}\\\textbf{optimisation}};
\node[anchor=east,align=right] at (-0.15,0.0) {\textbf{Learned}\\\textbf{policy}};

% ---- top lane: optimisation ----
\draw[ax] (0.1,2.1) -- (9.9,2.1);
\foreach \x/\d in {3.4/4.5, 5.6/6.7, 7.8/8.9}{
  \node[solve] at (\x,2.1) {solve};
  \node[dec] at (\d,2.1) {};
}
\draw[decorate,decoration={brace,amplitude=4pt},gray] (2.7,2.75) -- (8.5,2.75)
  node[midway,above=3pt,font=\scriptsize,gray]{cost paid at \emph{every} decision};

% ---- bottom lane: learned ----
\node[train] (tr) at (1.35,0.0) {Training\\(offline)};
\draw[ax] (2.6,0.0) -- (9.9,0.0);
\foreach \x/\d in {3.4/4.2, 5.2/6.0, 7.0/7.8, 8.6/9.2}{
  \node[fpass] at (\x,0.0) {fwd\\pass};
  \node[dec] at (\d,0.0) {};
}
\node[font=\scriptsize,gray,anchor=north] at (1.35,-0.85) {cost paid \emph{once}};
\node[font=\scriptsize,gray,anchor=north] at (6.4,-0.85) {each decision: one forward pass};
\end{tikzpicture}

    \caption[Where the cost of adapting to change is paid]{Where the cost of adapting to change is paid. A re-solved optimisation (top) recomputes the problem instance before every decision, so the cost recurs each time the channel or topology changes. A learned policy (bottom) pays the cost once, in an offline training phase, and reduces every runtime decision to a single forward pass through the trained model.}
    \label{fig:cost-shift}
\end{figure}

\subsection{Reinforcement Learning and Deep Q-Learning}
\label{subsec:reinforcement-learning}

\ac{RL} frames a decision problem as a sequence of interactions between an agent and an environment: at each step the agent observes the environment's state, selects an action, and receives a reward, with the goal of learning a policy that maximises the reward accumulated over time rather than at any single step~\cite{DBLP:books/lib/SuttonB18}.
%
For an offloading or placement decision, the state is a description of the current topology and workload, the action is a placement or offloading choice, and the reward encodes the objective the mapping should optimise -- latency, energy, or any of the criteria \cref{subsec:constraint-based-placement,subsec:energy-aware-orchestration} folded into an optimisation objective instead (\cref{fig:rl-offloading-loop}).

\begin{figure}[t]
    \centering
    % Reinforcement-learning agent-environment loop specialised to offloading.
% Requires: \usepackage{tikz}  \usetikzlibrary{arrows.meta,positioning}
\begin{tikzpicture}[
  >=Latex, font=\small,
  box/.style={draw,rounded corners,minimum width=34mm,minimum height=16mm,align=center,fill=gray!10},
]
\node[box] (agent) at (0,0) {\textbf{Agent}\\\scriptsize learned offloading policy};
\node[box] (env)   at (7.2,0) {\textbf{Environment}\\\scriptsize continuum topology + workload};

\draw[->,thick] (agent.north) to[bend left=26]
  node[above,align=center]{\emph{action}\\\scriptsize placement / offload choice} (env.north);
\draw[->,thick] (env.south) to[bend left=26]
  node[below,align=center]{\emph{state} \scriptsize (topology, load) \;/\; \emph{reward} \scriptsize ($-$latency, $-$energy)} (agent.south);
\end{tikzpicture}

    \caption[The reinforcement-learning loop for offloading]{The reinforcement-learning interaction loop, specialised to the offloading decision. The agent observes a state -- the current topology and workload -- selects an action -- a placement or offloading choice -- and receives a reward encoding the objective to optimise, such as negative latency or energy; over time it learns a policy that maximises the reward accumulated across the sequence of decisions.}
    \label{fig:rl-offloading-loop}
\end{figure}

%
Q-learning solves this problem by estimating, for every state-action pair, the expected cumulative reward of taking that action and then following the best policy thereafter -- the pair's \emph{Q-value} -- and updating that estimate from experience, without requiring a model of how the environment evolves~\cite{DBLP:journals/ml/WatkinsD92}.
%
Once the Q-values have converged, the policy they induce is immediate: in each state, take the action with the highest estimated value.
%
The original algorithm stores one Q-value per state-action pair in a table, which is workable only while the state space stays small enough to enumerate;
a topology described by the position, load, and connectivity of every node in a continuum-scale deployment has no such bound.
%
Deep Q-learning replaces the table with a neural network trained to approximate the Q-value function directly from raw or lightly processed state observations, letting the agent generalise to states it has not seen during training instead of memorising one value per state~\cite{DBLP:journals/nature/MnihKSRVBGRFOPB15}.
%
Training that network with the plain Q-learning update is unstable, because successive observations are strongly correlated and the update's own target drifts as the network changes.
%
Deep Q-learning restores stability with two mechanisms: an experience-replay buffer that samples past transitions at random to break the correlation, and a separate target network held fixed between periodic updates so the target stops chasing the estimate~\cite{DBLP:journals/nature/MnihKSRVBGRFOPB15}.
%
Mnih et al. demonstrate the resulting agent on Atari games played from raw pixel input, using the same network architecture and hyperparameters across a range of games with no problem-specific tuning, a generality that makes the algorithm exportable to a domain -- offloading over a continuum topology -- the original work never targeted~\cite{DBLP:journals/nature/MnihKSRVBGRFOPB15}.
%
Exporting it, however, still requires the state fed to the network to expose the structure of the problem: a topology is a graph of nodes and links, not a fixed-size vector of pixels, and the representation this difference calls for is the subject of \cref{subsec:gnns-for-edge}.

\subsection{Graph Neural Networks (GNNs) for Edge Topologies}
\label{subsec:gnns-for-edge}

A \ac{GNN} extends the neural architecture behind deep Q-learning to inputs that are graphs rather than fixed-size vectors, computing a representation for each node by repeatedly aggregating information from its neighbours~\cite{DBLP:conf/iclr/KipfW17}.
%
Because the aggregation is defined locally -- over a node and its neighbourhood rather than the graph as a whole -- one set of learned parameters applies to graphs that differ in size and connectivity, unlike an ordinary network whose input layer fixes the dimension of what it accepts.
%
This locality has made \acp{GNN} a standard tool for problems posed over communication-network topologies, where the graph is a physical arrangement of devices and links rather than an abstract relation~\cite{DBLP:journals/wc/LeeYDL22}.
%
Kipf and Welling formulate this aggregation as a graph convolution, propagating each node's features to its neighbours through an operation derived from a first-order approximation of spectral graph filtering, so that a node's learned representation comes to depend on the features of the nodes around it and, through further layers, on increasingly distant ones~\cite{DBLP:conf/iclr/KipfW17}.
%
Veli\v{c}kovi\'{c} et al. replace the fixed aggregation weights of the graph-convolutional formulation with an attention mechanism that learns how much weight one node's representation should give to each neighbour, rather than fixing that weight from the graph's connectivity alone~\cite{DBLP:conf/iclr/VelickovicCCRLB18} --
a property that matters once the nodes being aggregated over are not interchangeable, which they are not towards the edge of the continuum (\cref{subsec:continuum-characteristics}).
%
Both formulations, though, still assume a graph with a single type of node and a single type of link, an assumption a continuum topology does not meet: gateways, edge servers, and end devices differ in the resources and capabilities catalogued in \cref{subsec:continuum-characteristics}, and the links between them differ correspondingly in bandwidth and volatility.
%
Wang et al. address this mismatch with a \emph{heterogeneous} graph attention network (\cref{fig:gnn-aggregation}).
%
It learns both how much a node should attend to each of its neighbours and how much to weigh each type of relation that connects it to them,
rather than treating every neighbour and every link as an instance of the same relation~\cite{DBLP:conf/www/WangJSWCYY19}.
%
Applied to an edge-cloud topology, this two-level attention gives a placement or offloading policy a representation that mirrors the topology's own heterogeneity:
it distinguishes a link to a battery-limited robot from a link to a mains-powered gateway, rather than folding both into an undifferentiated edge~\cite{DBLP:conf/www/WangJSWCYY19}.

\begin{figure}[t]
    \centering
    % GNN neighbourhood aggregation: homogeneous vs. heterogeneous.
% Requires: \usepackage{tikz}
%   \usetikzlibrary{arrows.meta,positioning,fit,backgrounds}
\begin{tikzpicture}[
  >=Latex, font=\small,
  ctr/.style={circle,draw,thick,minimum size=8mm,inner sep=0pt,fill=white},
  nb/.style={circle,draw,minimum size=6.5mm,inner sep=0pt,fill=gray!15},
  cloud/.style={circle,draw,minimum size=7mm,inner sep=0pt,fill=gray!8},
  gw/.style={circle,draw,minimum size=7mm,inner sep=0pt,fill=gray!24},
  dev/.style={circle,draw,minimum size=7mm,inner sep=0pt,fill=gray!38},
  agg/.style={->,gray!70},
]
% ===== LEFT: homogeneous =====
\node[ctr] (v1) at (0,0) {$v$};
\node[nb] (a1) at (-1.5,1.1) {};
\node[nb] (a2) at ( 1.5,1.1) {};
\node[nb] (a3) at (-1.6,-1.0){};
\node[nb] (a4) at ( 1.6,-1.0){};
\foreach \n in {a1,a2,a3,a4}{ \draw[agg] (\n) -- (v1); }
\node[font=\scriptsize,align=center] at (-1.35,0.15) {$W$};
\node[anchor=north,align=center,font=\scriptsize,text width=3.6cm] at (0,-1.6)
  {\textbf{Homogeneous}\\one shared weight over identical neighbours};

% ===== RIGHT: heterogeneous =====
\begin{scope}[xshift=7.2cm]
\node[ctr] (v2) at (0,0) {$v$};
\node[cloud] (c) at (-1.6,1.1) {\scriptsize C};
\node[gw]    (g) at ( 1.6,1.1) {\scriptsize G};
\node[dev]   (d) at (-1.6,-1.0){\scriptsize D};
\node[dev]   (r) at ( 1.6,-1.0){\scriptsize R};
\draw[agg] (c) -- (v2);
\draw[agg] (g) -- (v2);
\draw[agg] (d) -- (v2);
\draw[agg] (r) -- (v2);
\node[font=\scriptsize] at (-0.6, 0.8)  {$\alpha_{\mathrm{cloud}}$};
\node[font=\scriptsize] at ( 0.6, 0.8)  {$\alpha_{\mathrm{gw}}$};
\node[font=\scriptsize] at (-0.7,-0.75) {$\alpha_{\mathrm{dev}}$};
\node[font=\scriptsize] at ( 0.7,-0.75){$\alpha_{\mathrm{dev}}$};
\node[anchor=north,align=center,font=\scriptsize,text width=4.2cm] at (0,-1.6)
  {\textbf{Heterogeneous}\\per-type weights + node/type attention $\alpha$};
\end{scope}
\end{tikzpicture}

    \caption[Neighbourhood aggregation: homogeneous vs.\ heterogeneous GNN]{Neighbourhood aggregation in a graph neural network. A homogeneous GNN (left) computes a node's representation by applying one shared weight to interchangeable neighbours. A heterogeneous GNN (right) distinguishes neighbours by type -- cloud (C), gateway (G), device (D), robot (R) -- and learns both how much a node should attend to each neighbour and how much to weigh each type of relation, so the learned representation mirrors the topology's own heterogeneity.}
    \label{fig:gnn-aggregation}
\end{figure}

%
Heterogeneity is not the only property of a continuum topology a \ac{GNN} must accommodate: the graph presented at run time is rarely the one seen during training, since nodes join, leave, and move (\cref{subsec:network-volatility}).
%
The spectral convolution of Kipf and Welling is transductive, trained on a fixed graph and not, on its own, defined for nodes or topologies absent at training time~\cite{DBLP:conf/iclr/KipfW17}.
%
Hamilton et al. lift this restriction with an inductive formulation that learns the aggregation function itself rather than an embedding bound to one graph, so that the trained model applies unchanged to nodes -- and to whole topologies -- it never saw~\cite{DBLP:conf/nips/HamiltonYL17}.
%
An inductive, heterogeneity-aware representation of the topology lets the deep Q-learning agent of \cref{subsec:reinforcement-learning}, trained offline, act on a continuum that keeps changing under it.
%
This is the combination developed into a full offloading policy in \cref{chap:heterogeneous-gnn}.

Whatever the method, the object of every decision reviewed in this chapter is the same kind of unit:
a service, a container, a function, a task --
inherited from the architectural paradigms whose deployment limitations were established in \cref{chap:edge-cloud-continuum}.
%
The problem posed in \cref{sec:the-deployment-problem} is therefore only half answered.
%
The optimisation and orchestration methods of \cref{sec:edge-orchestration-sota} and the learned policies of \cref{sec:ml-for-topologies} say how a mapping is computed and revised.
%
None of the reviewed work supplies a model of what is mapped at a granularity finer than these units,
matched to the collective programs of \cref{chap:programming-models}.
%
Part II begins with that model:
the deployable units are defined in \cref{chap:pulverization},
the strategies of \cref{chap:deployment-and-reconfiguration} and the offloading policy of \cref{chap:heterogeneous-gnn} then decide where those units run,
and one such deployment is built on a self-organising robot team in \cref{chap:demonstrator}.
